%%%%%%%%%%%%%%%%%%%%%%%%%%%%%%%%%%%%%%%%%%
\section{Discussion}
\label{sec:discussion}

\subsection{Interpretation of the factorial effects}
\label{subsec:discussion_factorial}

The interpretation will distinguish the effect of linguistic granularity from the effect of rule semantics. In particular, the discussion will assess whether increased resolution produces diminishing returns and whether the rule ranking changes across unimodal, scalable multimodal, and fixed-dimensional problem classes.

\subsection{Global versus problem-dependent configurations}
\label{subsec:discussion_global}

The central question is whether one static $(n,R)$ configuration remains optimal across the benchmark. Global and subgroup rankings will be compared using rank concordance and the observed transfer of the best configuration between problem classes. A heterogeneous ranking would indicate that static fuzzy-control semantics are not universally transferable.

\subsection{Threats to validity and limitations}
\label{subsec:limitations}

The conclusions are bounded by the selected benchmark suite, the fixed triangular Ruspini partitions, the eight rule patterns, the population size, and the chosen PSO baseline. Overlap, shoulder asymmetry, alternative defuzzification methods, and adaptive rule selection are outside the present factorial design. These limitations define direct extensions rather than uncontrolled variations of the current study.
